% from ega4-21c-fr.tex
\oldpage[IV-4]{327}
On est ainsi ramené au cas où $S=\operatorname{Spec}(A)$,
$X=\operatorname{Spec}(B)$, où $A$ est un anneau \emph{local}; puisque $X$ est
normal et $f$ lisse, $S$ est normal (17.5.7), donc $A$ est \emph{intègre et
intégralement clos}. Considérons alors $A$ comme limite inductive de ses
\emph{sous-$\mathbb Z$-algèbres de type fini}; comme la clôture intégrale d'une telle
sous-$\mathbb Z$-algèbre est un sous-anneau de $A$ et est aussi une
$\mathbb Z$-algèbre de type fini (7.8.3, (ii), (iii) et (vi)), $A$ est limite inductive
filtrante des sous-$\mathbb Z$-algèbres de type fini $A_\lambda$ de $A$ qui sont des
anneaux \emph{intégralement clos}. En vertu de (1.8.4.2), il existe un indice
$\lambda$ et une $A_\lambda$-algèbre de type fini $B_\lambda$ tels que
\[
  B=B_\lambda\otimes_{A_\lambda}A
\]
à un $A$-isomorphisme près. Posons $S_\lambda=\operatorname{Spec}(A_\lambda)$,
$X_\lambda=\operatorname{Spec}(B_\lambda)$; si $p_\lambda:X\to X_\lambda$ est la
projection canonique, on peut en outre supposer (puisque $Y$ est
\emph{constructible}) que $Y=p_\lambda^{-1}(Y_\lambda)$, où $Y_\lambda$ est une
partie fermée de $X_\lambda$ (8.3.11). Soit $f_\lambda$ le morphisme
$X_\lambda\to S_\lambda$; en vertu de la transitivité des fibres et de (4.2.6),
$Y_\lambda$ ne contient aucune composante irréductible des fibres
$f_\lambda^{-1}(s_\lambda)$ pour aucun $s_\lambda\in S_\lambda$. Puisque $f$ est
lisse, on peut aussi supposer que $f_\lambda$ est lisse (17.7.8); enfin, l'image du
point générique $\eta$ de $S$ dans $S_\lambda$ est le point générique $\eta_\lambda$
de $S_\lambda$, et on a
\[
  \operatorname{codim}((Y_\lambda)_{\eta_\lambda},
  (X_\lambda)_{\eta_\lambda})
  =\operatorname{codim}(Y_\eta,X_\eta)\geq2
\]
en vertu de (6.1.4). On voit donc que $S_\lambda$, $X_\lambda$ et $Y_\lambda$
vérifient toutes les hypothèses de (ii), et en posant
$X_\mu=X_\lambda\times_{S_\lambda}S_\mu$ pour $\lambda\leq\mu$, il en est de même de
$S_\mu$, $X_\mu$ et $Y_\mu$. Montrons que si l'on prouve que le couple
$(X_\mu,Y_\mu)$ est parafactoriel pour tout $\mu\gg\lambda$, il en est de même de
$(X,Y)$. En effet, soient $U=X-Y$, $U_\mu=X_\mu-Y_\mu$, $j:U\to X$,
$j_\mu:U_\mu\to X_\mu$ les injections canoniques; la projection
$p_\mu:X\to X_\mu$ étant un morphisme affine, on a
\[
  j_*(\mathcal O_U)=p_\mu^*((j_\mu)_*(\mathcal O_{U_\mu}))
  \qquad(\mathbf{II}, 1.5.2),
\]
et par suite l'homomorphisme canonique
$\rho:\mathcal O_X\to j_*(\mathcal O_U)$ n'est autre que $p_\mu^*(\rho_\mu)$, où
$\rho_\mu:\mathcal O_{X_\mu}\to(j_\mu)_*(\mathcal O_{U_\mu})$ est l'homomorphisme
canonique; ce dernier étant par hypothèse bijectif, il en est de même de $\rho$.
D'autre part, pour tout $\mathcal O_U$-Module inversible $\mathscr L$, il existe un
$\mu\gg\lambda$ et un $\mathcal O_{U_\mu}$-Module inversible $\mathscr L_\mu$ tels
que $\mathscr L=p_\mu'^*(\mathscr L_\mu)$ (en désignant par
$p_\mu':U\to U_\mu$ la restriction de $p_\mu$) (8.5.2 et 8.5.5), $U_\lambda$ étant
quasi-compact puisque $X_\lambda$ est noethérien. Par hypothèse, il existe un
$\mathcal O_{X_\mu}$-Module inversible $\mathscr L'_\mu$ tel que
$\mathscr L'_\mu|U_\mu$ soit isomorphe à $\mathscr L_\mu$;
$\mathscr L'=p_\mu^*(\mathscr L'_\mu)$ est alors un $\mathcal O_X$-Module
inversible tel que $\mathscr L'|U$ soit isomorphe à $\mathscr L$, ce qui prouve notre
assertion.

On est ainsi ramené à prouver (ii) lorsque l'anneau $A$ est une
$\mathbb Z$-algèbre de type fini intégralement close; comme les anneaux locaux
$\mathcal O_{S,s}$ de $S$ sont alors des anneaux excellents intégralement clos, leurs
\emph{complétés} sont aussi intégralement clos (7.8.3, (ii), (iii) et (vii)). Pour
prouver que le couple $(X,Y)$ est parafactoriel, il suffit de montrer que pour tout
$x\in Y$, l'anneau $\mathcal O_{X,x}$ est parafactoriel (21.13.10). Soit $y$ un point
fermé de $Y_{f(x)}$, qui soit spécialisation de $x$ (5.1.11); si l'on pose
$s=f(x)=f(y)$, $\mathcal O_{S,s}$ a un
